\section{Background and Experimental Setup}
\label{sec:background}

\subsection{AMD GPU Architectures}
\label{sec:bg-amd}

AMD has established two distinct GPU architectures targeting
different markets. The Compute DNA (CDNA) line is purpose-built for
data-center HPC, emphasizing high computational throughput and
high-bandwidth memory; the evolution from CDNA (MI50, MI100)
through CDNA~2 (MI210, MI250X) to CDNA~3 (MI300X) brought
substantial improvements in compute density and HBM bandwidth, with
the current flagship MI300X providing 304 compute units, 192~GB of
HBM3, and over 5~TB/s of memory bandwidth. The Radeon DNA (RDNA)
line targets gaming and consumer applications but is also relevant
for cost-sensitive HPC deployments and for development workflows;
RDNA~3 (RX~7900~XT) is the reference part here. Table~\ref{tab:gpus}
summarizes the four AMD GPUs evaluated in this paper.

\begin{table}[t]
\centering
\caption{AMD GPU architectures evaluated in this paper.}
\label{tab:gpus}
\footnotesize
\begin{tabular}{lllrrr}
\toprule
\textbf{GPU} & \textbf{Arch.} & \textbf{ISA} & \textbf{CUs} & \textbf{Memory} & \textbf{BW (GB/s)} \\
\midrule
MI50       & GCN    & gfx906  &  60 & 32 GB HBM2  & 1{,}024 \\
MI210      & CDNA~2 & gfx90a  & 104 & 64 GB HBM2e & 1{,}638 \\
MI300X     & CDNA~3 & gfx942  & 304 & 192 GB HBM3 & 5{,}300 \\
RX 7900 XT & RDNA~3 & gfx1100 &  84 & 20 GB GDDR6 &   800  \\
\bottomrule
\end{tabular}
\end{table}

The four reference parts span an order of magnitude in compute unit
count and memory bandwidth, and three distinct architectural
families. They share the SIMT execution model but differ on the
specific axes that compression kernels exercise the most: wavefront
size (CDNA fixed at 64 threads; RDNA configurable wave32 or wave64),
on-package cache capacity (4~MB on MI50 to 256~MB Infinity Cache on
MI300X), and host interconnect (PCIe Gen3 to Gen5 across the
lineup). These differences propagate directly into the
optimization choices presented in Sections~\ref{sec:optim-chunk-size}
and~\ref{sec:optim-adaptive-tiled}.

\subsection{Batched Chunked Compression on GPUs}
\label{sec:bg-batched}

The dominant GPU compression model, introduced by \nvcomp{} and
inherited by \arcto, is \emph{batched chunked} processing. The
input buffer is partitioned into fixed-size chunks (typically
64~KB) and the chunks are processed independently in parallel,
with one or more waves per chunk according to the algorithm's
internal structure. LZ4 uses a single wave per chunk with a
private hash table for match finding; Snappy uses two
cooperating waves; Cascaded applies a multi-stage pipeline
(delta plus run-length plus bit-packing) with 128 threads per
4~KB partition. The chunk size is the central parameter exposed
to the caller and is the primary lever for occupancy tuning.

At the API boundary, the caller hands the library a vector of
chunk pointers and an output device buffer. Two transfer
patterns are common: \emph{scattered pageable}, where each chunk
is staged from a separate pageable host allocation via its own
\texttt{hipMemcpyAsync}, and \emph{coalesced pinned}, where the
chunks are first packed into a single pinned host buffer and the
entire batch is uploaded in one bulk \texttt{hipMemcpyAsync}. The
pinned variant achieves PCIe peak bandwidth on the bulk copy but
pays a non-trivial allocation cost up front. The trade-off
between these two patterns is the subject of
Section~\ref{sec:optim-adaptive-tiled}.

\subsection{The Tooling Gap on AMD}
\label{sec:bg-gap}

The AMD software ecosystem centers on ROCm~\cite{rocm}, an
open-source platform that includes the HIP programming model. HIP source code is near-syntactically identical
to CUDA and can be retargeted to NVIDIA via a thin translation
header. However, while HIP simplifies porting CUDA code,
achieving optimal performance on AMD hardware often requires
architecture-specific optimizations that go beyond simple API
translation.

The compression-tool asymmetry between the two ecosystems is
documented in a recent characterization study on AMD GPU
architectures~\cite{anon_iccsa2026}: while NVIDIA has \nvcomp{}
as a production-ready library, the only AMD-side counterpart is
\hipcomp{}, a community-driven port started in 2022 that covers
only LZ4 and Snappy, applies direct source translation, and is
not maintained as a stable API. Scientists working on
AMD-equipped systems including El Capitan, Frontier, and others
cannot directly benefit from the mature CUDA compression
ecosystem and must either port code themselves, use suboptimal
tools, or forgo GPU-accelerated compression entirely. This
situation motivates the present work.

\subsection{Experimental Setup}
\label{sec:bg-setup}

All measurements in this paper use the same Singularity container
(\texttt{arcto\_gfx<arch>.sif}) built per AMD architecture from a
common Ubuntu~22.04 + ROCm~7.0.1 base. \arcto{} is built from
commit \texttt{f8b1c5f} of the \texttt{feature/scaling-instrumentation}
branch. The benchmark binaries time each phase of the
host-to-device path (allocation, pageable-to-pinned memcpy, bulk
H2D, compress kernel, decompress kernel) via host-side
\texttt{std::chrono} and GPU-side \texttt{hipEvent}, exposed by a
\texttt{--report\_phases} flag.

\paragraph{Hardware platforms}
The MI50, MI210, and MI300X experiments use a publicly accessible
HPC testbed at the Luxembourg site (vianden cluster), which
provides exclusive access to single-tenant compute nodes. The
RX~7900~XT experiments run on a workstation-class host at a host
institution research laboratory with a Threadripper PRO 5965WX
processor, 256~GB DDR4, and PCIe~Gen4 x16 host interconnect. Both
platforms run the same container image, differing only in the
architecture-specific build target.

\paragraph{Workloads}
The reference workload is the TTI (Tilted Transversely Isotropic)
seismic wavefield produced by the Fletcher forward-modelling
benchmark~\cite{fletcher2009tti}, the same workload used in the
prior cross-AMD characterization on which this paper
builds~\cite{anon_iccsa2026}. The dataset is
extracted from the middle of a $448^3$ grid simulation with 201
timesteps (to avoid the near-zero values dominant in
initialization timesteps that would inflate compression ratios)
and truncated to six canonical sizes covering the scaling range
of interest: \textbf{10~MB, 100~MB, 1~GB, 4~GB, 8~GB, 16~GB}. Not
every (GPU, size) combination is viable: the smaller-VRAM parts
cannot accommodate the largest inputs because the device input
plus output buffers and the compression scratch space together
exceed the on-package memory. Table~\ref{tab:viability} reports
the (GPU, size) combinations evaluated.

\begin{table}[t]
\centering
\caption{(GPU, size) combinations evaluated in this paper.
\checkmark{} = evaluated; ${\sim}$ = tight (input + output + scratch
near VRAM ceiling, evaluated with caution); --- = exceeds VRAM.
The benchmark allocates a peak of approximately $2 \times$ the input
size on the device (input buffer plus worst-case compressed-output
buffer, both held simultaneously during the compress kernel), so
the maximum input size that fits on a part of VRAM capacity $V$ is
about $V/2$.}
\label{tab:viability}
\footnotesize
\begin{tabular}{l c c c c c c}
\toprule
\textbf{GPU} & \textbf{10 MB} & \textbf{100 MB} & \textbf{1 GB} & \textbf{4 GB} & \textbf{8 GB} & \textbf{16 GB} \\
\midrule
MI50       (32 GB) & \checkmark & \checkmark & \checkmark & \checkmark & ${\sim}$ & --- \\
MI210      (64 GB) & \checkmark & \checkmark & \checkmark & \checkmark & \checkmark & ${\sim}$ \\
MI300X     (192 GB)& \checkmark & \checkmark & \checkmark & \checkmark & \checkmark & \checkmark \\
RX 7900 XT (20 GB) & \checkmark & \checkmark & \checkmark & \checkmark & --- & --- \\
\bottomrule
\end{tabular}
\end{table}

The RX~7900~XT was empirically confirmed to fail at 8~GB with
\texttt{hipErrorOutOfMemory}: the peak simultaneous device
allocation (input buffer at 8~GB plus the worst-case compressed
output buffer at $\approx 8.5$~GB plus the per-chunk scratch)
exceeds the 20~GB VRAM ceiling. The same physical constraint
applies to MI50 at 16~GB. The cross-architecture coverage of the
evaluation is therefore: small inputs validated on all four parts;
large inputs validated on the parts whose VRAM accommodates them.

\paragraph{Metrics and protocol}
Each (algorithm, size, mode) configuration runs in a separate
process invocation with 1 warmup iteration and 3--5 timed
iterations. The benchmark reports per-phase times in milliseconds
(allocation, pageable-to-pinned memcpy, bulk H2D, kernel),
end-to-end time as the sum of phases, and compression ratio.
Cross-mode comparisons use end-to-end speedup of the proposed mode
over the scattered-pageable baseline. The full schema is
documented in the benchmark CSV header and reproduced in the
public repository\footnote{
\TODO{public repo URL}}.
